\lysh{33826}{4}

\begin{alist}
\item Θέτοντας στην εξίσωση $x^4 - 8x^2 - 9 = 0$ όπου $x^2 = u \ge 0$, γίνεται:
$u^2 - 8u - 9 = 0$.
Το τριώνυμο έχει $\alpha = 1, \beta = -8, \gamma = -9$ και διακρίνουσα:
\[\Delta = \beta^2 - 4\alpha\gamma = (-8)^2 - 4 \cdot 1 \cdot (-9) = 64 + 36 = 100 > 0.\]
Οι ρίζες του είναι οι:
\[u_{1,2} = \dfrac{-\beta \pm \sqrt{\Delta}}{2\alpha} = \dfrac{-(-8) \pm \sqrt{100}}{2 \cdot 1} = \dfrac{8 \pm 10}{2} = \begin{cases} 9 > 0 \text{ δεκτή} \\ -1 < 0 \text{ απορρίπτεται} \end{cases}.\]

Όμως έχουμε $x^2 = u \Leftrightarrow x^2 = 9 \Leftrightarrow x = \pm \sqrt{9} \Leftrightarrow x = \pm 3$.

\item Όπως και στο ερώτημα $\alpha)$, η εξίσωση $x^4 - \beta x^2 + \gamma = 0$, αν θέσουμε όπου $x^2 = u$ με
$u \ge 0$, ισοδύναμα γίνεται:
$u^2 - \beta u + \gamma = 0$.

\begin{rlist}
\item Το τριώνυμο έχει διακρίνουσα $\Delta = \beta^2 - 4\gamma$, με:
$\beta^2 \ge 0$ και
$\gamma < 0$ οπότε $-4\gamma > 0$.
Συνεπώς, $\Delta > 0$ ως άθροισμα ενός μη αρνητικού και ενός θετικού αριθμού.

\item Επειδή $\Delta > 0$ το τριώνυμο έχει δύο άνισες ρίζες $u_1, u_2$. Από τους τύπους Vieta το
γινόμενο των ριζών είναι $P = \dfrac{\gamma}{\alpha} = \gamma < 0$. Άρα, οι ρίζες είναι μη μηδενικές και
ετερόσημες. Έστω $\begin{cases} u_1 < 0, \text{απορρίπτεται} \\ u_2 > 0 \text{ δεκτή} \end{cases}$.

Τότε έχουμε:
\[x^2 = u_2 \Leftrightarrow x = \pm \sqrt{u_2}.\]
\end{rlist}
\end{alist}

